\documentclass[11pt]{article}
\title{On the Closure of Compositional Hierarchies in Digital Symmetries}
\author{A rigorous researcher-engineer bridging abstract mathematics and concrete digital systems.}
\usepackage[margin=1in]{geometry}
\usepackage{graphicx}
\usepackage{amsmath,amssymb,mathtools}
\usepackage{tikz}
\usetikzlibrary{positioning,arrows.meta}
\usepackage{hyperref}
\graphicspath{{media/}{media/diagrams/}{images/}{artwork/}}

\begin{document}

\section*{10. Algebra of Programs \& Hardware DSLs}

This section develops the algebraic viewpoint that underlies program semantics and hardware domain-specific languages (DSLs): programs are not merely executed, but also interpreted as morphisms in an algebraic theory, transformed by equations, and validated by rewriting and property-based testing.

A useful starting point is the idea of a law coverage metric, written here as \(\varepsilon(r)\), which measures how well a test suite or mutation campaign exercises the intended equational laws of a representation \(r\). In the present setting, \(\varepsilon(r)\) can be read as a property-coverage or mutation-score proxy for laws: if a DSL is equipped with rewrite rules, normal forms, or algebraic identities, then a high value of \(\varepsilon(r)\) indicates that the generated tests are sensitive to violations of those laws. This is especially relevant for compiler passes and circuit optimizers, where a transformation may preserve extensional behavior while still breaking a stated equational invariant. The metric is therefore not a replacement for proof, but a practical complement to it: it helps quantify whether the laws that define the DSL are actually being exercised by the artifact.

The semantic backbone of many such DSLs is provided by Lawvere theories and their effectful extensions. A Lawvere theory presents an algebraic theory by specifying operations of finite arity together with equations, and it is well suited to describing compositional program fragments whose behavior is determined by generators and relations. When effects are present---for example, state, nondeterminism, exceptions, or resource usage---the pure equational picture must be enriched. In categorical terms, one often passes from ordinary algebraic theories to structures that account for effectful composition, while retaining the same guiding principle: syntax is generated freely, and semantics is obtained by quotienting by equations and interpreting in a model. This perspective connects naturally to operads and PROPs. Operads organize multi-input, single-output composition, while PROPs generalize to multiple inputs and multiple outputs, making them especially apt for circuits and signal-flow systems. Colored operads further refine the picture by tracking typed wires or heterogeneous interfaces. The common theme is that the algebraic presentation of a DSL should match the shape of the compositions it is meant to express.

For hardware DSL design, this algebraic viewpoint is not merely aesthetic; it is operational. A well-designed embedded DSL (EDSL) for circuits typically exposes a small set of combinators for primitive gates, wiring, duplication, deletion, and composition. The combinators should be chosen so that the resulting terms admit a useful normal form. Normal forms serve several purposes at once: they provide a canonical representation for equality checking, they support optimization by rewriting, and they make it possible to state and prove invariants about generated hardware. In practice, one often designs the DSL so that syntactic composition mirrors the intended circuit topology, while a separate normalization phase eliminates administrative overhead such as redundant identities, trivial permutations, or duplicated structure. The resulting artifact is easier to reason about because the algebra of the syntax aligns with the algebra of the semantics.

Rewriting systems supply the mechanism by which these normal forms are obtained. A rewrite rule is an oriented equation, and a collection of such rules induces a reduction relation on terms. Two classical properties govern the quality of a rewriting system. Termination ensures that repeated rewriting cannot continue indefinitely; confluence ensures that whenever a term can be rewritten in more than one way, all reduction paths can be joined to a common result. When both properties hold, every term has a unique normal form, and equational reasoning becomes algorithmic: two terms are equal modulo the theory precisely when their normal forms coincide. In the context of program and hardware DSLs, this is a powerful design target. It allows one to encode algebraic identities as rewrite rules, prove that the rules are sound with respect to the intended semantics, and then use normalization as a decision procedure for a large class of equivalences.

Completeness is the bridge between the abstract equational theory and the concrete rewrite system. An equational presentation is complete when the rewrite rules are sufficient to derive all valid equalities of the theory, typically modulo the chosen notion of equivalence. In a DSL for circuits, completeness means that the optimizer or normalizer is not merely performing ad hoc simplifications, but is systematically capturing the intended algebra. This is particularly important when the DSL is used as a compilation target or as a proof language for hardware transformations. If the rewrite system is incomplete, then some semantically valid transformations remain inaccessible to the toolchain; if it is unsound, then the toolchain may silently change meaning. The ideal is therefore a rewrite system that is both sound and complete for the fragment of the theory that the DSL is meant to represent, with termination and confluence providing the computational discipline needed for reliable automation.

An artifact-oriented implementation of these ideas can be organized around a small DSL with equational reasoning and QuickCheck-style randomized testing. The DSL exposes constructors for the primitive operations of the domain, together with an interpreter into a semantic model and a normalizer based on rewrite rules. Equational laws are then encoded as properties: for example, that normalization preserves denotation, that a rewrite step is semantics-preserving, or that two syntactically distinct expressions are observationally equivalent. QuickCheck-style generators can produce random well-typed terms, including terms near the boundaries of the rewrite system where bugs are most likely to appear. The property suite can then be scored using \(\varepsilon(r)\), giving a quantitative estimate of how thoroughly the laws are being exercised. In a mature workflow, the same laws appear in three places: as rewrite rules, as proof obligations, and as executable properties. Agreement among these three views is a strong indicator that the DSL is well founded.

The broader methodological lesson is that algebraic semantics, rewriting, and testing are not separate concerns but mutually reinforcing layers of the same design. Lawvere-theoretic syntax provides the generators and equations; operads and PROPs organize the shape of composition; rewriting turns equations into executable normalization; and property-based testing measures whether the intended laws are actually enforced by the implementation. For hardware DSLs in particular, this alignment is crucial because the cost of a semantic mismatch is not merely a failed test but a potentially incorrect circuit. The section therefore treats algebraic structure as an engineering resource: it constrains the language, enables optimization, and supports verification.

\paragraph{Artifact sketch.} A representative deliverable for this section is a small circuit DSL with the following components: a typed syntax of combinators; a denotational interpreter into Boolean or signal-flow semantics; a rewrite-based normalizer; a library of equational laws stated as executable properties; and a randomized test harness that reports law coverage via \(\varepsilon(r)\). Such an artifact makes the abstract discussion concrete and provides a reusable template for later sections on dataflow, reversible computation, and symmetry-constrained design.

\paragraph{Design note.} Relative to the surrounding chapters, this section serves as the bridge from symmetry and closure theory to executable language design. It intentionally emphasizes the operational role of equations: in the earlier material, algebraic structure classifies functions and circuits; here, the same structure becomes a programming discipline for building, normalizing, and testing DSLs. That bridge prepares the reader for later sections on dataflow graphs, reversible and quantum-flavored hierarchies, and algorithmic closure analysis.


\end{document}
